%% LyX trick_preamble_code_into_believing_that_this_was_created_by_lyx created this file.  For more info, see http://www.lyx.org/.
%% Do not edit unless you really know what you are doing.
\documentclass[a4paper,12pt]{article}
\usepackage[T1]{fontenc}
\usepackage[latin9]{inputenc}

\setlength{\parindent}{3mm}
\usepackage{setspace}
\onehalfspacing

\usepackage{amsmath}
\usepackage{url}
\usepackage{verbatim}
\usepackage[numbers]{natbib}
\usepackage{nomencl}
% the following is useful when we have the old nomencl.sty package
\providecommand{\printnomenclature}{\printglossary}
\providecommand{\makenomenclature}{\makeglossary}
\makenomenclature
\usepackage{varioref}
\usepackage{prettyref}
\usepackage{makeidx}
\makeindex

\usepackage{graphicx}
\usepackage{longtable}
\usepackage{xargs}

%%%%%%%%%%%%%%%%%%%%%%%%%%%%%% LyX specific LaTeX commands.
\providecommand{\LyX}{L\kern-.1667em\lower.25em\hbox{Y}\kern-.125emX\@}
\newcommand{\lyxline}[1][1pt]{%
  \par\noindent%
  \rule[.5ex]{\linewidth}{#1}\par}

\newcommand{\lyxarrow}{\leavevmode\,$\triangleright$\,\allowbreak}

%%%%%%%%%%%%%%%%%%%%%%%%%%%%%% User specified LaTeX commands.
\usepackage{subscript} % user specified as long as tex2lyx
% produces a format less than 408

\def\mycommand{\textquestiondown}

\begin{document}

\tableofcontents

\noindent This paragraph is not indented.

\section{References\index{References}}

Let's start with simple things: a label~\label{lab:test} and a
reference~\ref{lab:test}; note that I have tested ``unbreakable
space'' without warning (and worse than that, I just tested english
quotes too...).

Of course there are other kind of references, like page
reference~\pageref{lab:test}, but also equation
reference~\eqref{lab:test} (from amsmath package), or varioref's
equivalents~\vref{lab:test} and~\vpageref{lab:test}, without
forgetting pretty references like~\prettyref{lab:test}.

\section{Cites\index{Cites}}

Let's start with simple things: a \textbackslash{}cite: \cite[after]{article-crossref}
and two \textbackslash{}cites: \cite[after]{whole-set,article-crossref}

Now the natbib things:

\textbackslash{}citet: \citet[before][after]{article-crossref} \textbackslash{}citet{*}:
\citet*[before][after]{article-crossref} \textbackslash{}Citet: \Citet[before][after]{article-crossref}
\textbackslash{}Citet{*}: \Citet*[before][after]{article-crossref}

\textbackslash{}citep: \citep[before][after]{article-crossref} \textbackslash{}citep{*}:
\citep*[before][after]{article-crossref} \textbackslash{}Citep: \Citep[before][after]{article-crossref}
\textbackslash{}Citep{*}: \Citep*[before][after]{article-crossref}

\textbackslash{}citealt: \citealt[before][after]{article-crossref}
\textbackslash{}citealt{*}: \citealt*[before][after]{article-crossref}
\textbackslash{}Citealt: \Citealt[before][after]{article-crossref}
\textbackslash{}Citealt{*}: \Citealt*[before][after]{article-crossref}

\textbackslash{}citealp: \citealp[before][after]{article-crossref}
\textbackslash{}citealp{*}: \citealp*[before][after]{article-crossref}
\textbackslash{}Citealp: \Citep[before][after]{article-crossref}
\textbackslash{}Citealp{*}: \Citealp*[before][after]{article-crossref}

\textbackslash{}citeauthor: \citeauthor[before][after]{article-crossref}
\textbackslash{}citeauthor{*}: \citeauthor*[before][after]{article-crossref}
\textbackslash{}Citeauthor: \Citeauthor[before][after]{article-crossref}
\textbackslash{}Citeauthor{*}: \Citeauthor*[before][after]{article-crossref}

\textbackslash{}citeyear: \citeyear[before][after]{article-crossref}

\textbackslash{}citeyearpar: \citeyearpar[before][after]{article-crossref} 

\textbackslash{}nocite: \nocite{article-crossref}

\bibliographystyle{unsrt}
% Remove duplicate call of \bibliography since LaTeX throws an error.
% \bibliography{xampl}

With \textbackslash{}nocite\{{*}\}:
\bibliographystyle{unsrt}
\nocite{*}
\bibliography{xampl}

\section{Input files\index{Input files}}

We can input files too, like this 
\subsection{\label{sub:External-Subsection}External Subsection}

This is a small dummy child document to show how files can be inserted
to another document. Here are some accented characters to make sure
the encoding is passed to included files: ��
, or with the include
variant 
\subsection{\label{sub:External-Subsection}External Subsection}

This is a small dummy child document to show how files can be inserted
to another document. Here are some accented characters to make sure
the encoding is passed to included files: ��
 % unfortunately, including the doc twice
% generates a multiply defined label

If you prefer verbatim input, you can choose
between~\verbatiminput{foo} or~\verbatiminput*{foo}.

\section{URL\index{URL}\nomenclature[www]{URL}{uniform resource locator}}

An URL: \url{http://www.lyx.org}
\nomenclature{URL2}{uniform resource locator}

\section{Lists\index{Lists}}

\listoffigures
\listoftables

\section{Graphics\index{Graphics}}

There is also some basic support for graphics, in the form
\includegraphics{foo.eps}, or the slightly more elaborate
\includegraphics[bb=10bp 0bp 96bp 96bp,clip,height=1cm, width=1cm]{foo.eps}.

\section{Tables\index{Tables}}

The following example is stolen from the longtable documentation.
Since tex2lyx does not understand the special verbatim code that
was used in the original some lines have been rewritten using
\textbackslash textbackslash etc.
\let\package\textsf
\let\env\textsf
\let\code\texttt
\providecommand\finalclearpage{\clearpage}

\begin{longtable}{@{*}r||p{1in}@{*}}
KILLED & LINE!!!! \kill
\caption
[An optional table caption (used in the list of tables)]
{A long table\label{long}}\\
\hline\hline
\multicolumn{2}{@{*}c@{*}}%
     {This part appears at the top of the table}\\
\textsc{First}&\textsc{Second}\\
\hline\hline
\endfirsthead
\caption[]{(continued)}\\
\hline\hline
\multicolumn{2}{@{*}c@{*}}%
      {This part appears at the top of every other page}\\
\textbf{First}&\textbf{Second}\\
\hline\hline
\endhead
\hline
This goes at the&bottom.\\
\hline
\endfoot
\hline
These lines will&appear\\
in place of the & usual foot\\
at the end& of the table\\
\hline
\endlastfoot
\env{longtable}  columns  are specified& in the \\
same way as  in the \env{tabular}& environment.\\
\code{@\{*\}r||p\{1in\}@\{*\}}& in this case.\\
Each row ends with a& \code{\textbackslash\textbackslash} command.\\
The \code{\textbackslash\textbackslash} command  has an& optional\\
argument, just as in& the\\
 \env{tabular}&environment.\\[10pt]
See the effect of \code{\textbackslash\textbackslash[10pt]}&?\\
Lots of lines& like this.\\
Lots of lines& like this.\\
Lots of lines& like this.\\
Lots of lines& like this.\\
Also \code{\textbackslash hline} may be used,&  as in \env{tabular}.\\
\hline
That  was a \code{\textbackslash hline}&.\\
\hline\hline
That  was \code{\textbackslash hline\textbackslash hline}&.\\
\multicolumn{2}{||c||}%
{This is a \code{\textbackslash multicolumn\{2\}\{||c||\}}}\\
If a  page break  occurs at a \code{\textbackslash hline} then& a line is drawn\\
at the bottom of one  page  and at the& top of the next.\\
\hline
The \code{[t] [b] [c]} argument of \env{tabular}& can  not be used.\\
The optional argument may be  one of& \code{[l] [r] [c]}\\
to specify whether  the  table  should be& adjusted\\
to the  left, right& or centrally.\\
\hline\hline
Lots of lines& like this.\\
Lots of lines& like this.\\
Lots of lines& like this.\\
Lots of lines& like this.\\
Lots of lines& like this.\\
Lots of lines& like this.\\
Lots of lines& like this.\\
Lots of lines& like this.\\
Lots of lines& like this.\\
Lots of lines& like this.\\
Lots of lines& like this.\\
Lots of lines& like this.\\
Lots of lines& like this.\\
Lots of lines& like this.\\
Lots of lines& like this.\\
Lots of lines& like this.\\
Lots of lines& like this.\\
Lots of lines& like this.\\
Lots of lines& like this.\\
Lots of lines& like this.\\
Some lines may take up a lot of space, like this: &
    \raggedleft This last column is a ``p'' column so this
    ``row'' of the table can take up several lines. Note however that
    \TeX\ will  never break a page within such a row. Page breaks only
    occur between rows of the table or at \code{\textbackslash hline} commands.
    \tabularnewline
Lots of lines& like this.\\
Lots of lines& like this.\\
Lots of lines& like this.\\
Lots of lines& like this.\\
Lots of lines& like this.\\
Lots of lines& like this.\\
Lots of lines& like this.\\
\hline
Lots\footnote{This is a footnote.} of lines& like this.\\
Lots   of   lines& like this\footnote{\env{longtable} takes special
    precautions, so that footnotes may also be used in `p' columns.}\\
\hline
Lots of lines& like this.\\
Lots of lines& like this.
\end{longtable}

\section{Macros}

LyX supports several kinds of macros:
def \def\macroa#1{a #1 a}
global def \global\def\macrob#1{b #1 b}
long def \long\def\macroc#1{c #1 c}
global long def \global\long\def\macrod#1{d #1 d}
providecommand \providecommand{\macroe}[1]{e #1 e}
providecommand* \providecommand*{\macrof}[1]{f #1 f}
newcommand \newcommand{\macrog}[1]{g #1 g}
renewcommand \renewcommand{\macrog}[1]{h #1 h}
newcommand* \newcommand*{\macroi}[1]{i #1 i}
renewcommand* \renewcommand*{\macroi}[1]{j #1 j}
providecommandx \providecommandx{\macrok}[1]{k #1 k}
providecommandx* \providecommandx*{\macrok}[1]{l #1 l}
newcommandx \newcommandx{\macrom}[1]{m #1 m}
renewcommandx \renewcommandx{\macrom}[1]{n #1 n}
newcommandx* \newcommandx*{\macroo}[1]{o #1 o}
renewcommandx* \renewcommandx*{\macroo}[1]{p #1 p}
DeclareRobustCommand \DeclareRobustCommand{\macroq}[1]{q #1 q}
DeclareRobustCommand* \DeclareRobustCommand*{\macror}[1]{r #1 r}
DeclareRobustCommandx \DeclareRobustCommandx{\macros}[1]{s #1 s}
DeclareRobustCommandx* \DeclareRobustCommandx*{\macrot}[1]{t #1 t}

Now use them all:
\macroa{x} \macrob{x} \macroc{x} \macrod{x} \macroe{x} \macrof{x} \macrog{x}
\macroi{x} \macrok{x} \macrom{x} \macroo{x} \macroq{x} \macror{x}
The following tow don't work, but they should???
%\macros{x} \macrot{x}

\section{Special formattings\index{Special formattings}}

\subsection{LyX line}

test

{\tiny \lyxline{\tiny}}{\tiny \par}

test {\Huge Test} {\tiny test} test

\lyxline{\normalsize}

test

{\Huge \lyxline{\Huge}}{\Huge \par}

test

\subsection{Line breaks}

They can also or be broken by a newline\\
or by a starred newline \\*
or by a newline with space, comment and argument \\ %hu
[3cm]
or by a newline command \newline
or by a line break \linebreak
or by a defined line break \linebreak % again with a comment
	[4]

There are even newlines with weird arguments, but these are not
handled by LyX\\*[1cm]
so we try to use ERT in this case.

\subsection{Page breaks}

They can also or be broken by a newpage \newpage
or by a page break \pagebreak
or by a defined page break \pagebreak % again with a comment
	[4]

\section{Special characters\index{Special characters}}

Then one has those macros with a long name for a short meaning, like
\textasciitilde, \textasciicircum{} or \textbackslash{}, \slash{},
\nobreakdash- and the characters
that LaTeX wants to espace because they are active, like \_\&\#\$\{\}\%.

And what about special characters like hyphe\-nation mark,
ellipsis\ldots, and end-of-sentence\@. LyX also supports a menu
separator\lyxarrow{}and a spif\textcompwordmark{}fy ligature break.

LyX translates the phrases LyX, TeX, LaTeX2e and LaTeX
to the commands \LyX{}, \TeX{}, \LaTeXe{} and \LaTeX{}.
If these phrases occur as part of other words (like 1LyX or aTeX or LaTeX3)
they should not be put into ERT.

Test for whitespace handling of commands: The following lines should
result in identical output:

builtin \textasciicircum{} unicodesymbols \j{} user \mycommand{} xx\par
builtin \textasciicircum {} unicodesymbols \j {} user \mycommand{} xx\par
builtin \textasciicircum % with a comment
{} unicodesymbols \j % and a second one
{} user \mycommand % and another
{} xx

A sub\textsubscript{sc\emph{ript}} and super\textsuperscript{script
with $a^2+b^2=c^2$ math}.

\printindex{}

\printnomenclature{}

\end{document}
